% ==========================================================
% ENTREGA 1 — A pergunta do trabalho
% Teoria do Aprendizado Estatístico — Ciências de Dados
% Fatec Baixada Santista – Rubens Lara
% ==========================================================

\documentclass[
  article,12pt,oneside,a4paper,english,brazil
]{abntex2}

\usepackage[T1]{fontenc}
\usepackage[utf8]{inputenc}
\usepackage{lmodern}
\usepackage{amsmath,amssymb}
\usepackage{booktabs}
\usepackage{graphicx}
\usepackage[alf]{abntex2cite}
\usepackage{natbib}
\hypersetup{hidelinks}

\titulo{\textit{É viável prever o volume mensal de comércio exterior (em kg) a partir de município, fluxo, mês, país e produto?}}
\tituloestrangeiro{}
\autor{Equipe FFV}
\local{Santos-SP}
\data{\today}

\newcommand{\nomecurso}{Tecnologia em Ciência de Dados}
\newcommand{\nomedisciplina}{Teoria do Aprendizado Estatístico}
\newcommand{\nomeprofessor}{Prof.\ Dr. João Paulo Ferreira de Mello}
\newcommand{\repositorio}{TAE-FFV}
\newcommand{\integrantes}{%
  Felipe Steinfeld -- RA 0051352511017 \\
  Flávia Alessandra da Silva Barbosa -- RA 0051352511018 \\
  Victor Gabriel Leme da Silva -- RA 0051352421012%
}

\begin{document}



\selectlanguage{brazil}
\frenchspacing
\maketitle

\begin{center}
  \small
  Fatec Baixada Santista -- Rubens Lara \\
  \nomecurso{} \quad|\quad \nomedisciplina{} \quad|\quad \nomeprofessor{} \\[0.4em]
  \integrantes \\[0.4em]
  Repositório: \texttt{https://github.com/flaalevia15-gif/TAE-FFV/tree/main/estrutura/notebooks}
\end{center}

\textual

\section{Introdução}

O comércio exterior constitui uma dimensão relevante da atividade econômica
dos municípios relacionados à região do Porto de Santos, envolvendo fluxos de
importação e exportação associados a diferentes países e produtos. Neste
trabalho são analisados dados dos municípios de Cubatão, Guarujá e Santos,
considerando o município de domicílio fiscal das empresas declarantes. Essa
distinção é importante: os dados municipais do Comex Stat representam o
domicílio fiscal do exportador ou importador e não necessariamente o local
físico onde a mercadoria foi produzida ou movimentada
\citep{mdic_comexstat_faq}.

A base utilizada é proveniente do Comex Stat, sistema oficial de divulgação
das estatísticas de comércio exterior de bens do Brasil, mantido pelo
Ministério do Desenvolvimento, Indústria, Comércio e Serviços (MDIC). O
sistema permite consultas por município, país, produto, período e quilograma
líquido, além da exportação dos resultados em CSV
\cite{mdic_comexstat,mdic_dados_abertos}. Essas características tornam a base
adequada para investigar se informações disponíveis antes da observação do
volume podem auxiliar sua previsão.

A proposta também se relaciona aos ODS 8 e 9, associados ao crescimento
econômico, produtividade, infraestrutura e inovação
\cite{onu_ods8,onu_ods9}. Assim, este trabalho investiga se é viável prever o
volume mensal de comércio exterior, em quilogramas, a partir de município,
fluxo, mês, país e produto.

\section{O banco de dados e o problema}

A base principal desta entrega é a \texttt{df\_sazonal\_lideres.csv}, derivada
dos dados do Comex Stat. Ela contém registros mensais dos principais pares de
município, fluxo, país e produto selecionados na preparação dos dados. O
período observado vai de janeiro de 2019 a julho de 2026, totalizando
481 observações. Cada linha representa uma combinação mensal de
município, operação de comércio exterior, país e produto, e a variável
{Quilograma Líquido} representa o volume mensal correspondente.

As principais variáveis são \textbf{Município}, \textbf{Fluxo},
\textbf{mes\_numero}, \textbf{País}, \textbf{Código SH4} e
\textbf{Quilograma Líquido}. A descrição textual do SH4 e o
{Valor US\$ FOB} também estão disponíveis na base, mas não fazem parte
dos preditores definidos para esta pergunta. A descrição do SH4 é redundante
com seu código e o valor FOB é uma medida contemporânea do mesmo registro,
podendo fornecer informação que não estaria disponível na formulação original
da pergunta.

A escolha desta base ocorre porque ela preserva a dimensão temporal mensal e,
ao mesmo tempo, mantém as características de município, fluxo, país e
produto que motivaram a pergunta. A existência de 481 observações para cinco
preditores candidatos fornece quantidade suficiente para uma primeira etapa
de aprendizado supervisionado, embora a validação deva respeitar a ordem
temporal.

O problema prático é apoiar análises de comércio exterior e planejamento
logístico. Um analista ou gestor poderia utilizar uma estimativa do volume
mensal esperado para antecipar demanda operacional, comparar períodos e
identificar combinações de rotas com maior movimentação. A previsão não deve
ser interpretada como previsão direta da movimentação física do Porto de
Santos, pois a base municipal utiliza o domicílio fiscal das empresas
declarantes \cite{mdic_comexstat_faq}.

\section{A pergunta}

A pergunta: \textit{é possível prever o volume mensal
de comércio exterior, em quilogramas, para uma combinação de município,
fluxo, país e produto, usando o mês como informação de previsão?}

\begin{table}[htb]
  \centering
  \caption{Ficha da pergunta}
  \label{tab:ficha}
  \small
  \begin{tabular}{@{}p{0.30\textwidth}p{0.62\textwidth}@{}}
    \toprule
    \textbf{Campo} & \textbf{Preenchimento} \\
    \midrule
    1. A pergunta, em uma frase &
    É possível prever o volume mensal de comércio exterior em kg a partir de
    município, fluxo, mês, país e produto? \\
    2. Resposta $Y$ &
    \texttt{Quilograma Líquido}; quantitativa; kg. \\
    3. Preditores $X$ &
    \texttt{Município}, \texttt{Fluxo}, \texttt{mes\_numero},
    \texttt{País} e \texttt{Código SH4}; $p=5$. \\
    4. Tipo &
    Aprendizado supervisionado; regressão; predição. \\
    5. Métrica &
    RMSE, em quilogramas. \\
    6. Linha de base &
    Prever sempre a média: 106.622.215 kg \\
    7. Dono do problema &
    Analista/gestor de comércio exterior ou logística que precise estimar
    volumes mensais para planejamento e comparação de operações. \\
    \bottomrule
  \end{tabular}
  \fonte{Elaborada pelos autores.}
\end{table}

A escolha por regressão decorre do fato de que a resposta é uma quantidade
numérica contínua, e não uma classe. O objetivo é \emph{predição}, pois a
finalidade é produzir valores estimados para novas observações, e não medir
efeitos causais ou testar uma hipótese sobre um coeficiente específico.
Como os dados possuem dimensão temporal, uma etapa posterior deverá avaliar o
modelo com separação temporal, evitando que informações de meses futuros
participem do treinamento de uma previsão de meses anteriores.

\section{Defesa da pergunta}

A Tabela~\ref{tab:diag} resume as verificações iniciais. O conjunto possui
481 observações e cinco preditores definidos para a pergunta. Não há
valores faltantes nas variáveis utilizadas.

\begin{table}[htb]
  \centering
  \caption{Diagnóstico do banco}
  \label{tab:diag}
  \small
  \begin{tabular}{@{}p{0.32\textwidth}p{0.60\textwidth}@{}}
    \toprule
    \textbf{Verificação} & \textbf{Resultado} \\
    \midrule
    Tamanho ($n \times p$) &
    481 $\times$ 5 \\
    Linha de base &
    RMSE = 106.622.215 kg \\
    Vazamento &
    Nenhum dos cinco preditores contém diretamente a resposta. \\
    Níveis dos preditores &
    Município: 3; Fluxo: 2; Mês: 12; País: 4; Código SH4: 4 \\
    Valores faltantes &
    Nenhum \\
    \bottomrule
  \end{tabular}
  \fonte{Elaborada pelos autores.}
\end{table}

\textbf{Vazamento.} Não foi identificado vazamento direto entre os cinco
preditores escolhidos e a resposta. Entretanto, {Valor US\$ FOB} foi
mantido fora de $X$, pois é uma informação do próprio registro comercial e
apresenta forte associação com o volume. Utilizá-lo mudaria a pergunta
original e poderia tornar a avaliação excessivamente dependente de uma
informação contemporânea. Da mesma forma, {Descrição SH4} não foi
usada porque duplica a informação do código do produto.

\textbf{Tamanho.} São 481 observações para cinco preditores. Esse número é
adequado para uma primeira investigação, mas os preditores são em grande
parte categóricos e possuem poucos níveis. Portanto, a quantidade de linhas
não deve ser interpretada como 481 situações completamente independentes:
existem apenas três municípios, dois fluxos, quatro países e quatro produtos
nos principais pares selecionados. A validação deve considerar essa
estrutura.

\textbf{Balanço e faltantes.} O balanço de classes não se aplica, pois a
resposta é quantitativa. Também não foram encontrados valores ausentes nas
variáveis utilizadas, portanto não é necessário imputar dados nesta etapa.

A distribuição da resposta é bastante assimétrica: a média é de
85.409.831 kg
kg, enquanto a mediana é de
50.288.228 kg
kg, e o maior valor chega a
756.741.873 kg
kg. A diferença entre média e mediana indica a presença de meses com volumes
muito superiores aos valores típicos.

\begin{figure}[htb]
  \centering
  \caption{Distribuição do volume mensal de comércio exterior}
  \label{fig:resposta}
\includegraphics[width=0.70\linewidth]{fig-resposta.png}
  \fonte{Elaborada pelos autores.}
\end{figure}

A transformação logarítmica é usada apenas para visualizar melhor a
distribuição, pois os volumes variam em várias ordens de grandeza. O
diagnóstico sustenta a pergunta, mas indica que modelos futuros deverão
considerar a forte assimetria da resposta e, principalmente, a natureza
temporal dos dados.

\textbf{Veredito.} A pergunta se sustenta como problema de regressão
supervisionada: há uma resposta quantitativa bem definida, cinco preditores
disponíveis e ausência de valores faltantes. A principal cautela é
metodológica: como o objetivo é prever volumes mensais, os testes futuros
devem treinar com períodos anteriores e avaliar em períodos posteriores.
Além disso, a base representa os municípios pelo domicílio fiscal das
empresas, de modo que as conclusões devem ser apresentadas como previsões do
comércio exterior associado aos municípios, e não como medição direta da
movimentação física do porto. A pergunta é interessante pois aplica conhecimentos de aprendizado
estatístico e faz parte da área de ciência de dados pois um analista pode usar
a predição para ter uma ideia antecipada sobre o comércio.

\postextual
\bibliography{referencias}

\end{document}
